\section{Effective synthesis of an optimal interior learner}
\label{sec:effectivelearning80}

The finite readout in Theorem~\ref{thm:interiorlearnedcode79} does not
by itself give computable entries for its measurement operators.
We now construct a finite algebraic readout, and then a finite
trusted-control realization, at the joint query and description
orders.  The construction uses exact quantifier elimination.  The
tomography estimator of \cite{MB67} supplies an existence bound;
we do not compile its continuous outcome measurement or assume
that its coarse-grained operators have computable entries.

\subsection{A decidable collective-readout problem}
Fix integers $d,m\ge1$.  Prepare $m$ independent maximally
entangled probe--reference pairs and apply the unknown measurement
to each probe.  Put $E_1=E$ and $E_0=I_d-E$.  The complete
classical--quantum record is
\begin{equation}\label{eq:choirecord80}
 \Gamma_m(E)=\bigoplus_{y\in\{0,1\}^m}R_y(E),\qquad
 R_y(E)=d^{-m}\bigotimes_{\ell=1}^m E_{y_\ell}^{\mathsf T}.
\end{equation}
The transpose is in the public input basis.  The blocks are
\emph{subnormalized}: their traces are the probabilities of the
classical strings $y$.  The real and imaginary parts of every entry
of $R_y(E)$ are polynomials with rational coefficients in the
$d^2$ real coordinates of $E$.  No division by an unknown outcome
probability occurs.

A final $L$-outcome readout consists of a POVM
$(A_{j,y})_{j=1}^L$ on the $d^m$ retained reference dimensions,
selected by the observed classical string $y$.  Its outcome
probabilities are
\begin{equation}\label{eq:readoutpolynomial80}
 p_j(E;A)=\sum_y\operatorname{tr}(A_{j,y}R_y(E)),\qquad
 A_{j,y}\succeq0,\quad\sum_{j=1}^L A_{j,y}=I_{d^m}.
\end{equation}
This description uses the outcome bits classically.  It entails no
coherent access to those bits or to a measurement environment.
Conversely, any collective POVM on \eqref{eq:choirecord80} can be
dephased in $y$ without changing its probabilities, giving precisely
such a family.

\begin{lemma}[Effective algebraic readout selection]
\label{lem:algebraicreadout80}
Given $d,m$, finitely many Gaussian-rational Hermitian centres
$C_1,\ldots,C_L$, and rational $a>0$, $0<\alpha<1$, it is
decidable whether a readout \eqref{eq:readoutpolynomial80} satisfies
\begin{equation}\label{eq:readoutfeasibility80}
 \sum_{j:\,\|E-C_j\|_{\rm op}\le a}p_j(E;A)\ge1-\alpha
 \qquad\hbox{for every }E\in\mathfrak I_d.
\end{equation}
When it is feasible, one finite deterministic computation returns
such a family with algebraic entries, valid for every real
$E\in\mathfrak I_d$.
\end{lemma}
\begin{proof}
Parameterize each Hermitian matrix by its real coordinates.
Positivity is a finite system of polynomial inequalities: use all
principal minors, or the equivalent real symmetric representation
for a complex Hermitian matrix.  Trace and tensor products in
\eqref{eq:readoutpolynomial80} yield real polynomials with rational
coefficients.  The same is true of the input constraint
$I_d/4\preceq E\preceq3I_d/4$ and of
\[
 \mathsf{Good}_j(E):\quad
 aI_d-(E-C_j)\succeq0,\qquad
 aI_d+(E-C_j)\succeq0.
\]
For each subset $S\subseteq\{1,\ldots,L\}$ impose the implication
\[
 \left(\bigwedge_{j\in S}\mathsf{Good}_j(E)\right)
 \wedge\left(\bigwedge_{j\notin S}\neg\mathsf{Good}_j(E)\right)
 \quad\Longrightarrow\quad
 \sum_{j\in S}p_j(E;A)\ge1-\alpha.
\]
The conjunction of these finitely many implications, universally
quantified over $E\in\mathfrak I_d$, is exactly
\eqref{eq:readoutfeasibility80}.  Equality belongs to
$\mathsf{Good}_j$; its complement uses the corresponding strict
inequalities.  Hence boundary cases require no numerical tolerance.

Eliminate the $E$ variables while retaining the readout entries as
free variables, and include the positivity and normalization
conditions in \eqref{eq:readoutpolynomial80}.  Effective real
quantifier elimination decides whether this rational semialgebraic
feasible set is empty.  If it is nonempty, the real-closed-field
transfer principle gives a point with real algebraic coordinates.
One explicit selection rule uses a fixed effective enumeration of
algebraic tuples and tests the resulting quantifier-free formula
by exact algebraic sign determination.  It terminates on this
nonempty set.
These are standard effective real-algebraic operations; see
\cite[Theorems~2.77, 2.80 and Chapter~14]{BPR80}.
The eliminated formula is equivalent over the reals, so the selected
point satisfies the condition for all real effects, not merely
algebraic ones.  The unknown effect is a quantified variable, never
supplied data or an oracle queried by this computation.
\end{proof}

\subsection{Finite controls at the joint optimal rates}
\begin{theorem}[Effective optimal interior learning and description]
\label{thm:effectiveinterior80}
Let $d,N\ge1$ be integers, and let
$0<\delta\le2^{-13}$ and $0<\eta\le1/8$ be rational.
There is a common learner on $\mathfrak I_d$ whose entire classical
specification is computed from $(d,N,\delta,\eta)$ in finitely
many steps.  In the finite trusted-control model of
Section~\ref{sec:finitecontrol78}, it returns a fixed-length word
decoded by a public deterministic decoder into a legal rational
effect $C$, with
\begin{align}
 \sup_{E\in\mathfrak I_d}
  \mathbb P_E\{d_N(E,C)>\delta\}&\le\eta,
                                         \label{eq:effectiverisk80}\\
 M&\le C_0\frac{N}{\delta^2}
                \left(d^2+d\log\frac1\eta\right),
                                         \label{eq:effectivecalls80}\\
 B&=\frac{d^2}{2}\log_2N+d^2\log_2(1/\delta)+O(d^2).
                                         \label{eq:effectivepayload80}
\end{align}
All constants are absolute.  The number $M$ and the finite readout
are fixed before any unknown-device call.  Acquisition uses fresh
one-call probe--reference pairs; only completed outputs undergo
collective processing.  Thus the construction belongs to every
fresh-block class $1\le b\le N$.

The call order matches the coherent adaptive lower bound of
Theorem~\ref{thm:interiorconfidence80}, and the description order
matches Theorem~\ref{thm:interiorentropy79}.  No bound on classical
search time, algebraic-description length, quantum workspace, or
elementary known-gate count is asserted.
\end{theorem}
\begin{proof}
Construct the finite rational dictionary $\mathcal D=(C_1,\ldots,C_L)$
of Theorem~\ref{thm:interiorcodec79} at parameters $(d,N,\delta)$.
Set
\[
 k=\lceil\sqrt N\rceil,
 \qquad a=\frac{\delta}{8k},\qquad\alpha=\eta/2.
\]
For $m=1,2,\ldots$, apply Lemma~\ref{lem:algebraicreadout80} to
these centres, $a$, and $\alpha$.  At the first feasible $m$,
extract its algebraic readout and set $M=m$.  Each infeasible
candidate is rejected by a terminating decision procedure; there
is no unbounded search for a witness at such a candidate.

We prove termination and the call bound without computing any
operator of the imported tomography procedure.  Apply
Lemma~\ref{lem:binaryconfidenceupper80} at operator accuracy
$\epsilon=\delta/(64k)$ and failure $\eta/2$.  Its independent
Choi acquisition uses
\[
 m_0\le C_1 k^2\delta^{-2}
                 \left(d^2+d\log\frac2\eta\right)
\]
calls and returns a legal effect $F$.  Clip $F$ spectrally to
$[I_d/4,3I_d/4]$, obtaining $T$.  On the estimation success event,
Weyl's inequality gives $\|T-F\|_{\rm op}\le\epsilon$ and hence
$\|T-E\|_{\rm op}\le2\epsilon$.  The rounding and predecessor
selection in Theorem~\ref{thm:interiorcodec79} produce an index $j$
with $\|T-C_j\|_{\rm op}\le5\delta/(64k)$.  Therefore
\begin{equation}\label{eq:effectiveexistence80}
 \|E-C_j\|_{\rm op}
 \le\frac{7\delta}{64k}<\frac{\delta}{8k}=a
\end{equation}
on that event.

Clipping, coordinate rounding with its prescribed ties, and
selection from a finite list are Borel operations.  Composing them
with the final tomography measurement defines a finite-outcome
POVM on the $m_0$ Choi records.  Its dephased form is a readout
\eqref{eq:readoutpolynomial80} satisfying
\eqref{eq:readoutfeasibility80}.  Only its existence is used here;
its entries and any defining integrals are not inputs to the
algebraic search.  Thus the search stops with $M\le m_0$.
Since $k^2\le4N$ and
$\log(2/\eta)\le C_2\log(1/\eta)$, this proves
\eqref{eq:effectivecalls80}.

The selected ideal readout has probability at least $1-\eta/2$ of
returning a centre at operator distance at most $a$ from $E$.
The target lies in $\mathfrak I_d$ and every centre lies in
$[I_d/8,7I_d/8]$.  Lemma~\ref{lem:interiormetric78} therefore gives
\[
 d_N(E,C_j)\le4\sqrt N\,a\le\delta/2
\]
on this event.

It remains to realize the computed readout with finite trusted
controls.  For each classical string $y$, its algebraic POVM has
the algebraic isometry
\[
 V_y\psi=\sum_{j=1}^L|j\rangle_C|j\rangle_Z
                         \sqrt{A_{j,y}}\,\psi.
\]
Tracing out $Z$ and the factor space gives the required classical
readout.  This is a finite table, including strings of probability
zero under some targets.  Algebraic positive square roots and the
finite descriptions of Section~\ref{sec:finitecontrol78} make each
entry effectively approximable.  In particular, the normalization
of a rational approximation to an isometry preserves the exact
channel constraints; a fixed final dephasing preserves the
classical output type.

There are $M$ fresh-pair preparation events and one final readout
event.  Require total unhalved error at each event at most
\begin{equation}\label{eq:effectivecontrolbudget80}
 e_{\rm event}=\frac{\eta}{M+1}.
\end{equation}
For a preparation this is trace-norm error in the fresh joint pair;
for the readout it is complete diamond-norm error, uniformly over
all $y$ and all retained reference states.  Split each allowance
equally between approximation by a finite rational specification
and its actual trusted realization.  Any additional elementary
instructions are charged by the same total pathwise budget $\eta$.
Lemma~\ref{lem:controlhybrid78} then bounds the total variation
between ideal and actual output indices by $\eta/2$.
Applying it to the original bad-output event $d_N(E,C_j)>\delta$
proves \eqref{eq:effectiverisk80}.  No continuity of a decision
boundary or perturbation of the returned effect is assumed.

Finally transmit the dictionary index using
$B=\lceil\log_2L\rceil$ bits.  The exact public decoder, including
its legal values on unused words, is unchanged.
Theorem~\ref{thm:interiorcodec79} gives the upper order in
\eqref{eq:effectivepayload80}.  Since the dictionary covers
$\mathfrak I_d$, Theorem~\ref{thm:interiorentropy79} gives its
matching lower order.  More generally, positive success probability
at every target forces any fixed finite decoder's range to cover
the target family.  This proves the stated optimality claims.
\end{proof}

The synthesis is an effective construction of a new collective
learner at the existing tomography query order.  It is not a claim
of efficient compilation of the original tomography measurement.
At the selected $M$, there are $2^M$ possible classical strings,
each selecting an $L$-outcome measurement on dimension $d^M$.
The algebraic search, its coefficients, the corresponding isometry
descriptions, and their preparation work are separate resources
from the transmitted index.  The per-event norm tolerance in
\eqref{eq:effectivecontrolbudget80} introduces no additional
$\delta/N$ requirement; dimensions and coordinate precision still
depend on $\delta$ through $M$ and $L$.  Fresh preparation remains
tensor-separated from the stored references, so finite-control
approximation preserves the one-call acquisition restriction.
